% Generated draft: requires mathematical and source review.
\documentclass[11pt]{article}
\usepackage[margin=25mm]{geometry}
\usepackage{fontspec}
\setmainfont{Latin Modern Roman}
\usepackage{amsmath,amssymb,mathtools,unicode-math}
\setmathfont{Latin Modern Math}
\usepackage{enumitem,needspace,xurl,hyperref}
\setlength{\parindent}{0pt}
\setlength{\parskip}{5pt}
\setlength{\emergencystretch}{3em}
\newcommand{\N}{\mathbb{N}}
\newcommand{\Z}{\mathbb{Z}}
\newcommand{\Q}{\mathbb{Q}}
\newcommand{\R}{\mathbb{R}}
\newcommand{\C}{\mathbb{C}}
\newcommand{\F}{\mathbb{F}}
\newcommand{\A}{\mathbb{A}}
\newcommand{\PP}{\mathbb{P}}
\newcommand{\E}{\mathbb{E}}
\newcommand{\eps}{\varepsilon}
\newcommand{\dd}{\,\mathrm{d}}
\begin{document}
\section*{UnsolvedMath ID: 801134}
\textbf{Source problem code:} catalogue entry
\subsection{Definitions and mathematical statement}
Let \(C_5\) be the cycle graph on vertices \(1,2,3,4,5\), with edges \(\{1,2\},\{2,3\},\{3,4\},\{4,5\},\{5,1\}\). Its right-angled Artin group is
\[
A(C_5)=\left\langle g_1,g_2,g_3,g_4,g_5\ \middle|\ [g_i,g_{i+1}]=e\ (1\leq i\leq 5)\right\rangle,
\]
where indices are read modulo \(5\) and \([a,b]=aba^{-1}b^{-1}\). No relations are imposed on pairs corresponding to nonedges.

The source asks: \emph{Is the pentagon right-angled Artin group \(\mathrm{II}_1\)-factor stable?} It places this question in a discussion of Jung-type properties but does not define the stability property in the supplied excerpt. Consequently, the following is an explicit conditional formulation, whose agreement with the intended source definition requires review.

Under the interpretation that \(\mathrm{II}_1\)-factor stability means uniform tracial stability of the defining unitary relations, let \(M\) range over all \(\mathrm{II}_1\) factors, with normalized trace \(\tau\), unitary group \(\mathcal U(M)\), and norm
\[
\|x\|_{2,\tau}=\tau(x^*x)^{1/2}.
\]
Does the following assertion hold?
\[
\begin{gathered}
\forall\varepsilon>0\ \exists\delta>0\ \forall (M,\tau)\ \forall u_1,\ldots,u_5\in\mathcal U(M),\\
\max_{1\leq i\leq5}\|u_i u_{i+1}-u_{i+1}u_i\|_{2,\tau}<\delta\\
\Longrightarrow\quad
\exists v_1,\ldots,v_5\in\mathcal U(M):\\
v_i v_{i+1}=v_{i+1}v_i\quad(1\leq i\leq5),\qquad
\max_{1\leq i\leq5}\|u_i-v_i\|_{2,\tau}<\varepsilon,
\end{gathered}
\]
with \(u_6=u_1\) and \(v_6=v_1\)? Here \(\delta\) must be independent of the factor and the unitary tuple. The exact tuple defines a group homomorphism \(A(C_5)\to\mathcal U(M)\); injectivity is not required by this conditional formulation. The excerpt supplies no separability restriction or further hypotheses, and does not establish that this interpretation is the intended one.
\subsection{Short English statement}
The source asks whether the right-angled Artin group on a pentagon is \(\mathrm{II}_1\)-factor stable. Under the explicitly conditional interpretation above, this asks whether five unitaries whose adjacent pairs almost commute can always be perturbed, uniformly over all \(\mathrm{II}_1\) factors, to five unitaries whose adjacent pairs commute exactly. The intended definition of stability must be checked against the source.
\subsection{Sources}
\begin{itemize}
\item Local UnsolvedMath registry record: permanent ID \(801134\), exact source problem code \texttt{catalogue entry}; title: \emph{Ultraproduct II\_1 Factors --- Problem 35}. The record attributes the problem to Elayavalli, with proposed year \(2025\), and supplies the status label \texttt{open}. This label is source metadata, not an independently verified mathematical assessment.
\item Supplied source list: Elayavalli, \emph{50 Open Problems: Ultraproduct II\_1 Factors} (2025), Problem 35, \url{https://arxiv.org/abs/2511.20377}. The record reports access on 2026-10-16 and extraction by \texttt{arxiv\_latex}; these are supplied provenance statements, not a claim of independent source access.
\item The source statement refers to \texttt{ScottSri2019ultraproduct} and \texttt{AGKE} for detailed discussion of the relevant Jung-type properties and what is known and unknown, and also cites \texttt{JungTubularity}. These reference labels are preserved exactly. Their bibliographic entries and links are not supplied, so their identities and precise definitions require verification.
\end{itemize}
\end{document}
